\documentclass[11pt,letterpaper]{article}
\usepackage[margin=1in]{geometry}
\usepackage{graphicx} % Required for inserting images
\usepackage{tikz}
\usepackage{amsthm,amsmath,amssymb}
\usepackage{algorithm}
\usepackage{algpseudocode}
\usepackage{xcolor}
\usepackage{soul}
\usepackage{enumitem}

\usepackage[hidelinks]{hyperref}
\usepackage{cleveref}

% Theorems
\newtheorem{theorem}{Theorem}[section]
\newtheorem{corollary}[theorem]{Corollary}
\newtheorem{lemma}[theorem]{Lemma}
\newtheorem{claim}[theorem]{Claim}
\theoremstyle{definition}
\newtheorem{definition}[theorem]{Definition}
\newtheorem{remark}[theorem]{Remark}
\newtheorem{observation}[theorem]{Observation}

\crefname{theorem}{Theorem}{Theorems}
\Crefname{theorem}{Theorem}{Theorems}
\crefname{corollary}{Corollary}{Corollaries}
\Crefname{corollary}{Corollary}{Corollaries}
\crefname{lemma}{Lemma}{Lemmas}
\Crefname{lemma}{Lemma}{Lemmas}
\crefname{claim}{Claim}{Claims}
\Crefname{claim}{Claim}{Claims}
\crefname{definition}{Definition}{Definitions}
\Crefname{definition}{Definition}{Definitions}
\crefname{remark}{Remark}{Remarks}
\Crefname{remark}{Remark}{Remarks}
\crefname{observation}{Observation}{Observations}
\Crefname{observation}{Observation}{Observations}

\newcommand{\reals}{\mathbb{R}}
\newcommand{\Span}{\mathrm{span}}
\newcommand{\coreq}{\mathrm{coreq}}

\newcommand{\calA}{\mathcal{A}}
\newcommand{\calB}{\mathcal{B}}
\newcommand{\calC}{\mathcal{C}}
\newcommand{\calD}{\mathcal{D}}
\newcommand{\calE}{\mathcal{E}}
\newcommand{\calF}{\mathcal{F}}
\newcommand{\calG}{\mathcal{G}}
\newcommand{\calH}{\mathcal{H}}
\newcommand{\calI}{\mathcal{I}}
\newcommand{\calJ}{\mathcal{J}}
\newcommand{\calK}{\mathcal{K}}
\newcommand{\calL}{\mathcal{L}}
\newcommand{\calM}{\mathcal{M}}
\newcommand{\calN}{\mathcal{N}}
\newcommand{\calO}{\mathcal{O}}
\newcommand{\calP}{\mathcal{P}}
\newcommand{\calQ}{\mathcal{Q}}
\newcommand{\calR}{\mathcal{R}}
\newcommand{\calS}{\mathcal{S}}
\newcommand{\calT}{\mathcal{T}}
\newcommand{\calU}{\mathcal{U}}
\newcommand{\calV}{\mathcal{V}}
\newcommand{\calW}{\mathcal{W}}
\newcommand{\calX}{\mathcal{X}}
\newcommand{\calY}{\mathcal{Y}}
\newcommand{\calZ}{\mathcal{Z}}

\newcommand{\cost}{\mathrm{COST}}
\newcommand{\ALG}{\mathrm{ALG}}
\newcommand{\OPT}{\mathrm{OPT}}
\newcommand{\Flow}{\textsc{Flow}}
\newcommand{\ouralg}{\textsc{LineTracker}}

\newcommand{\EE}{\mathbb{E}}

\newcommand{\interval}[1]{\lfloor #1 \rceil}
\newcommand{\dis}[1]{#1}
\newcommand{\off}{\mathrm{off}}
\newcommand{\on}{\mathrm{on}}
\newcommand{\org}{\mathrm{org}}
\newcommand{\rel}{\mathrm{rel}}
\newcommand{\new}{\mathrm{new}}

\newcommand{\solution}{\medskip\noindent\textbf{Solution.}\par\medskip}

\newcommand{\Var}{\mathrm{Var}}

\title{21-701: Discrete Mathematics \\ Homework 4}
\author{Poon Tze-Yang}
\date{Due November 2, 2026}


\begin{document}

\maketitle

\noindent Instructor: Wesley Pegden

\begin{enumerate}[series=q]
\item Let $A$ be a set of $n$ real numbers, and write
\[
A + A \cdot A = \{a + bc \mid a, b, c \in A\}.
\]
Prove that $|A + A \cdot A| \geq cn^{3/2}$ for some constant $c > 0$.
\end{enumerate}

\solution

Let $S=A+A\cdot A$, and let $d=|S|$. 
Consider the $A\times S$ grid, which has $nd$ points.
Let the set of such points be $P$
For each $a,b\in A$, define the line $y = ax+b$.
Let the set of such lines be $\calL$, so there are $n^2$ such lines.
For each $c\in A$ and each line $ax+b$, the line is incident to the point $(c,ax+b)$ in the grid.
Hence, the number of line-point incidences is $\geq n^3$.
Szemer\'edi-Trotter upper bounds the number of line-point incidences, so
\begin{align*}
    n^3\leq k(|\calL|^\frac{2}{3}|P|^\frac{2}{3}+|\calL|+|P|)
\end{align*}
for sufficiently large constant $k$.
This simplifies to:
\begin{align*}
    n^3\leq k(n^2d^\frac{2}{3} + nd + n^2),
\end{align*}
which implies $d\geq kn^\frac{3}{2}$ for sufficiently large $k>0$.


\begin{enumerate}[resume=q]
\item Prove that in a drawing of a graph $G$ with the smallest possible number of crossings, no two edges cross more than once.
\end{enumerate}

\solution

Suppose we have a drawing of a graph where the edges $\gamma_1,\gamma_2: [0,1]\to \reals^2$ cross at points $p=\gamma_1(s_1) = \gamma_2(s_2)$ and $q=\gamma_1(t_1) = \gamma_2(t_2)$, for some $0<s_1<t_1<0$ and $0<s_2<t_2<0$.
We replace the $(s_1,t_1)$ portion of $\gamma_1$ with the $(s_2,t_2)$ portion of $\gamma_2$, and vice versa, to obtain $\gamma_1'$ and $\gamma_2'$.
Assume no three edges cross at a point, so for some small disc $D$ around $p$, the only edges in $D$ are $\gamma_1'$ and $\gamma_2'$.
The two edges only intersect at the point $p$, so they can be redrawn within $D$ to not intersect.
A similar change can be done at $q$.
Then the total number of crossings has decreased by two, and no new crossings are created.
Hence, this process can be repeated until there are no pairs of edges that cross more than once.


\begin{enumerate}[resume=q]
\item Let $M$ be a $(2 \times n)$-matrix with nonzero real entries such that any two columns of $M$ form a $(2 \times 2)$-matrix of nonzero determinant. Show that $M$ has at most $O(n^{4/3})$ pairs of columns that form a $(2 \times 2)$-matrix of determinant $1$.

\emph{Hint:} Think of the columns as points in $\mathbb{R}^2$. Now transform this into an incidence problem.
\end{enumerate}

\solution
We will instead upper bound the pairs of columns that form a matrix of determinant $\pm 1$, which suffices for the claim.
Let each column be a point in $\reals^2$. 
Let this set of $n$ points be $P$.
For each column $(a b)^\top$, if column $(x y)^\top$ pairs with it to form a matrix of determinant $\pm 1$, then
\begin{align*}
    ay-bx = \pm 1 \Rightarrow y = \frac{b}{a}x \pm \frac{1}{a},
\end{align*}
so there are two such lines for each point in $P$.
Let this set of lines be $L$, where $|L|=2n$.

Every point-line incidence corresponds to a pair of columns whose determinant is $\pm1$.
Let the number of incidences be $I(L,P)$.
By Szemer\'edi-Trotter,
\begin{align*}
    I(L,P) &\leq k(|\calL|^\frac{2}{3}|P|^\frac{2}{3}+|\calL|+|P|)\\
        &=O(n^\frac{4}{3}).
\end{align*}

\begin{enumerate}[resume=q]
\item Let $\calF$ be a set of $n$ lines in $\mathbb{R}^3$, such that no four lines have a common point of intersection, and no three lines lie in a common plane. Show that there is a subset $\calF' \subset \calF$ of lines such that no three of them have a common point of intersection and with cardinality $|\calF'| \geq cn^{3/4}$ for some absolute constant $c > 0$.

\emph{Hint:} Consider a random subset of $\calF$.
\end{enumerate}

\solution


\begin{enumerate}[resume=q]
\item For this problem you show an \emph{upper} bound on the size of a smallest Kakeya set in $\mathbb{F}_q^n$. In particular, assuming $q$ is odd (although a similar construction can be made when $q$ is even), define
\begin{gather*}
I = \{(a_1,\dots,a_{n-1},0) \in \mathbb{F}_q^n\},\\
J = \{(a_1,\dots,a_{n-1},b) \in \mathbb{F}_q^n \mid \forall i = 1,\dots,n-1,\ a_i + b^2 \text{ is a square in } \mathbb{F}_q\}
\end{gather*}
and let
\[
K = I \cup J.
\]
\begin{enumerate}[label=\arabic*.]
\item Prove that for every direction $\mathbf{b} \in \mathbb{F}_q^n \setminus \mathbf{0}$, $K$ contains a line of the form $\{\mathbf{a} + \mathbf{b}t \mid t \in \mathbb{F}_q\}$ for some $\mathbf{a}$ ($K$ is a Kakeya set in $\mathbb{F}_q^n$). (Hints on last page.)
\item Prove that $|K| = \frac{1}{2^{n-1}} q^n + O(q^{n-1})$.
\end{enumerate}
\end{enumerate}

\solution


\newpage

Hints for part 1 of Question 5: Note: for directions $\mathbf{b}$ whose last coordinate is $0$, you can just use the set $I$ to make your line. For other directions you will use $J$. $q$ is odd means the characteristic of the field is bigger than $2$. The \emph{characteristic} of the field is the minimum number $j$ such that adding $x \neq 0$ to itself $j$ times always gives $0$---this is the same for any element $x$ of the field. The characteristic of $\mathbb{F}_q$ for $q = p^k$ is $p$. In particular, when the characteristic of $\mathbb{F}_q$ is bigger than $2$, for any element $0 \neq c \in \mathbb{F}_q$ there is an element $0 \neq 2c = c + c \in \mathbb{F}_q$ (note that $2$ doesn't refer to an element of our field but to repeated addition).

When choosing what $\mathbf{a}$ to start your line at, you might as well choose the last coordinate $a_n = 0$, and when choosing the others in such a way that all the points $\mathbf{a} + \mathbf{b}t$ lie in $J$, you can use division by the number $2b_n$, which may be useful given the definition of $J$.

\end{document}
